\section{Discussion}
\label{sec:discussion}

\subsection{Pourquoi l'optimisation des caractéristiques aide}

La recherche de séparabilité en Phase~1 sélectionne les paramètres de descripteurs
qui maximisent l'information discriminante entre classes avant tout entraînement
de classifieur. Cela évite un essai-erreur coûteux en Phase~2 et garantit que
les caractéristiques entrant dans le benchmark sont déjà bien réglées. Le gain de
\FusionGain{} points de pourcentage du HOG seul au Full(opt) --- modeste en valeur
absolue --- est statistiquement significatif (McNemar vs.\ RF : $p = \McNemarRF$)
et suggère que les quatre familles de descripteurs complémentaires capturent une
information non redondante.

\subsection{Pourquoi le SVM performe le mieux}

Le SVM à noyau polynomial capture des frontières de classe non linéaires dans
l'espace Full(opt) à 526 dimensions. Les descripteurs denses en haute dimension
(le HOG contribue à la majorité des 526 dimensions) favorisent généralement les
apprenants à marge plutôt que les stratégies de bagging ou de boosting sur des
jeux de données de taille modeste. Le fait que le KNN figure également en tête
du classement suggère que l'espace Full(opt) possède une géométrie locale forte
près des frontières de décision.

\subsection{Limites}

\begin{itemize}
  \item Le jeu de données (\NTotal{} images) reste modeste ; les résultats peuvent
        ne pas se généraliser à des populations multi-scanners ou multi-centres.
  \item Le pipeline de prétraitement (CLAHE, redimensionnement fixe) n'est pas adapté
        aux protocoles d'acquisition spécifiques.
  \item Les modèles d'apprentissage profond dépassent souvent le ML classique sur
        de grands jeux IRM ; notre pipeline cible le scénario de données limitées
        et de contraintes d'interprétabilité.
\end{itemize}

\subsection{Implications pratiques}

Le pipeline s'exécute entièrement sur CPU, sans outillage propriétaire, et produit
des descripteurs dont la signification physique (texture, fréquence, contours) est
compréhensible par les experts du domaine. Il est donc adapté comme outil d'aide
à la décision dans des contextes où les ressources de calcul ou les volumes de
données excluent l'apprentissage profond.
